\textbf{COLUNI 2023 -- Questão 07}\par\medskip Na introdução do livro “Histórias que os jornais não contam” (2017), Moacyr Scliar, definiu suas crônicas publicadas no jornal Folha de São Paulo como “histórias que esqueceram de acontecer”. Considerando as características próprias do gênero em questão, produzidas pelo autor, na obra acima indicada, é INCORRETO afirmar que:\begin{enumerate}
\item o gênero crônica está situado entre o universo jornalístico e o literário. Cumpre função jornalística, se o
      autor desejar, assim como apresenta uma natureza literária, já que o cronista recria o fato cotidiano por
      meio da leveza, da beleza, da poesia, da crítica, do humor.
\item por sua exclusividade argumentativa, o gênero crônica pode provocar uma reflexão em relação a um fato
      do cotidiano, marcado pela peculiaridade do autor no que se refere ao seu tom humorístico, irônico e
      aparentemente despretensioso.
\item assim como outros gêneros, a crônica é uma construção que está intimamente relacionada ao contexto
      em que é produzida e, portanto, as escolhas sobre seu tema, estilo e construção são guiadas pela
      dinâmica social, abordando temas ligados à vida cotidiana.
\item por sua natureza reflexiva, por organizar-se em textos/enunciados curtos, breves, por lidar com fatos do
      cotidiano, e utilizar uma linguagem simples e coloquial, na crônica, o cronista se aproxima do leitor.\end{enumerate}
